% GENERATED FILE. Edit canonical lessons, then run npm run generate:books.
% canonical-source-hash: fnv1a64:021a8cfe58c06f5e
% canonical-lessons: SA-C233-varsati, SA-C233-rohati, SA-C233-pratiksate, SA-C233-anugacchati, SA-C233-nirgacchati

\chapter{Rains, Sprouts, Waits, Follows, Goes out}
\label{ch:sa-rains-sprouts-waits-follows-goes-out}
\hlchaptermodality{233}

\begin{chapteropening}
\textbf{By the end of this chapter:} \emph{I can say rains, sprouts, waits, follows and goes out.}

Complete the last lesson of chapter 233: I can say rains, sprouts, waits, follows and goes out.
\end{chapteropening}

\section[\texorpdfstring{\sk{वर्षति} (varṣati)}{वर्षति (varṣati)}]{\sk{वर्षति} (varṣati) — rains}

\label{lesson:SA-C233-varsati}



\begin{quote}
Before the new one: say the Sanskrit for is lost, then the Sanskrit for gets angry.
\end{quote}

\subsection*{You'll want to know: \sk{वर्षति}}
\textbf{\sk{वर्षति}} — \emph{varṣati} — \textquotedblleft{}rains\textquotedblright{}.

\begin{cousinweb}
\emph{varṣati} is the present of the root √vṛṣ, \textquotedblleft{}to rain\textquotedblright{} (MW), a Rigvedic verb (CDIAL 11394). It lives on as Hindi \emph{barasnā}, \textquotedblleft{}to rain\textquotedblright{}, and its noun \emph{varṣa}, \textquotedblleft{}rain, year\textquotedblright{}, gives Hindi \emph{baras}, \textquotedblleft{}year\textquotedblright{} (CDIAL 11392).
\end{cousinweb}

\begin{grammarlens}[title={What you've built}]
One more word for this chapter.
\end{grammarlens}

\subsection*{Your turn}
\begin{itemize}
\raggedright
  \item \emph{Say it:} \emph{varṣati}
  \item \emph{Say it:} \emph{varṣati}, once more
  \item \emph{Say it:} say \emph{varṣati}
  \item \emph{Recall:} say the Sanskrit for is lost, then the Sanskrit for gets angry, then say \emph{varṣati} again
\end{itemize}

\begin{tcolorbox}[breakable,colback=teal!4,colframe=teal!35!black,title={Before you move on}]
What does \sk{वर्षति} mean? (\textquotedblleft{}Rains\textquotedblright{}.) Say it once more.
\end{tcolorbox}
\section[\texorpdfstring{\sk{रोहति} (rohati)}{रोहति (rohati)}]{\sk{रोहति} (rohati) — sprouts, grows}

\label{lesson:SA-C233-rohati}



\begin{quote}
Before the new one: say the Sanskrit for gets angry, then the Sanskrit for rains.
\end{quote}

\subsection*{You'll want to know: \sk{रोहति}}
\textbf{\sk{रोहति}} — \emph{rohati} — \textquotedblleft{}sprouts, grows\textquotedblright{}.

\begin{cousinweb}
\emph{rohati} is the present of the root √ruh, \textquotedblleft{}to rise, grow, mount\textquotedblright{} (MW). In the Rigveda it means both \textquotedblleft{}grows\textquotedblright{} and \textquotedblleft{}climbs\textquotedblright{} (CDIAL 10862), and Hindi \emph{rohanā} means \textquotedblleft{}to mount, climb\textquotedblright{}.
\end{cousinweb}

\begin{grammarlens}[title={What you've built}]
One more word for this chapter.
\end{grammarlens}

\subsection*{Your turn}
\begin{itemize}
\raggedright
  \item \emph{Say it:} \emph{rohati}
  \item \emph{Say it:} \emph{rohati}, once more
  \item \emph{Say it:} say \emph{rohati}
  \item \emph{Recall:} say the Sanskrit for gets angry, then the Sanskrit for rains, then say \emph{rohati} again
\end{itemize}

\begin{tcolorbox}[breakable,colback=teal!4,colframe=teal!35!black,title={Before you move on}]
What does \sk{रोहति} mean? (\textquotedblleft{}Sprouts, grows\textquotedblright{}.) Say it once more.
\end{tcolorbox}
\section[\texorpdfstring{\sk{प्रतीक्षते} (pratīkṣate)}{प्रतीक्षते (pratīkṣate)}]{\sk{प्रतीक्षते} (pratīkṣate) — waits}

\label{lesson:SA-C233-pratiksate}



\begin{quote}
Before the new one: say the Sanskrit for rains, then the Sanskrit for sprouts.
\end{quote}

\subsection*{You'll want to know: \sk{प्रतीक्षते}}
\textbf{\sk{प्रतीक्षते}} — \emph{pratīkṣate} — \textquotedblleft{}waits\textquotedblright{}.

\begin{cousinweb}
\emph{pratīkṣate} is built on \emph{prati}, \textquotedblleft{}towards\textquotedblright{}, and the root √īkṣ, \textquotedblleft{}to look\textquotedblright{} (MW): to wait is to look out for someone. It means \textquotedblleft{}beholds\textquotedblright{} in the Atharvaveda and \textquotedblleft{}waits for\textquotedblright{} in the Taittirīya Saṃhitā (CDIAL 8620); Marathi \emph{paikhṇẽ} means \textquotedblleft{}to wait for\textquotedblright{}.
\end{cousinweb}

\begin{grammarlens}[title={What you've built}]
One more word for this chapter.
\end{grammarlens}

\subsection*{Your turn}
\begin{itemize}
\raggedright
  \item \emph{Say it:} \emph{pratīkṣate}
  \item \emph{Say it:} \emph{pratīkṣate}, once more
  \item \emph{Say it:} say \emph{pratīkṣate}
  \item \emph{Recall:} say the Sanskrit for rains, then the Sanskrit for sprouts, then say \emph{pratīkṣate} again
\end{itemize}

\begin{tcolorbox}[breakable,colback=teal!4,colframe=teal!35!black,title={Before you move on}]
What does \sk{प्रतीक्षते} mean? (\textquotedblleft{}Waits\textquotedblright{}.) Say it once more.
\end{tcolorbox}
\section[\texorpdfstring{\sk{अनुगच्छति} (anugacchati)}{अनुगच्छति (anugacchati)}]{\sk{अनुगच्छति} (anugacchati) — follows}

\label{lesson:SA-C233-anugacchati}



\begin{quote}
Before the new one: say the Sanskrit for sprouts, then the Sanskrit for waits.
\end{quote}

\subsection*{You'll want to know: \sk{अनुगच्छति}}
\textbf{\sk{अनुगच्छति}} — \emph{anugacchati} — \textquotedblleft{}follows\textquotedblright{}.

\begin{cousinweb}
\emph{anugacchati} is built on \emph{anu}, \textquotedblleft{}after\textquotedblright{}, and \emph{gacchati}, \textquotedblleft{}goes\textquotedblright{}, from the root √gam (MW). To follow is to go after; √gam is the root of \emph{gacchāmi}, \textquotedblleft{}I go\textquotedblright{}, which you have already met.
\end{cousinweb}

\begin{grammarlens}[title={What you've built}]
One more word for this chapter.
\end{grammarlens}

\subsection*{Your turn}
\begin{itemize}
\raggedright
  \item \emph{Say it:} \emph{anugacchati}
  \item \emph{Say it:} \emph{anugacchati}, once more
  \item \emph{Say it:} say \emph{anugacchati}
  \item \emph{Recall:} say the Sanskrit for sprouts, then the Sanskrit for waits, then say \emph{anugacchati} again
\end{itemize}

\begin{tcolorbox}[breakable,colback=teal!4,colframe=teal!35!black,title={Before you move on}]
What does \sk{अनुगच्छति} mean? (\textquotedblleft{}Follows\textquotedblright{}.) Say it once more.
\end{tcolorbox}
\section[\texorpdfstring{\sk{निर्गच्छति} (nirgacchati)}{निर्गच्छति (nirgacchati)}]{\sk{निर्गच्छति} (nirgacchati) — goes out}

\label{lesson:SA-C233-nirgacchati}



\begin{quote}
Before the new one: say the Sanskrit for waits, then the Sanskrit for follows.
\end{quote}

\subsection*{You'll want to know: \sk{निर्गच्छति}}
\textbf{\sk{निर्गच्छति}} — \emph{nirgacchati} — \textquotedblleft{}goes out\textquotedblright{}.

\begin{cousinweb}
\emph{nirgacchati} is built on \emph{nis}, \textquotedblleft{}out\textquotedblright{}, and \emph{gacchati}, \textquotedblleft{}goes\textquotedblright{}, from the root √gam (MW): \textquotedblleft{}goes out, comes forth\textquotedblright{}. The \emph{s} of \emph{nis} becomes \emph{r} before \emph{g}.
\end{cousinweb}

\begin{grammarlens}[title={What you've built}]
One more word for this chapter.
\end{grammarlens}

\subsection*{Your turn}
\begin{itemize}
\raggedright
  \item \emph{Say it:} \emph{nirgacchati}
  \item \emph{Say it:} \emph{nirgacchati}, once more
  \item \emph{Say it:} say \emph{nirgacchati}
  \item \emph{Recall:} say the Sanskrit for waits, then the Sanskrit for follows, then say \emph{nirgacchati} again
\end{itemize}

\begin{tcolorbox}[breakable,colback=teal!4,colframe=teal!35!black,title={Before you move on}]
What does \sk{निर्गच्छति} mean? (\textquotedblleft{}Goes out\textquotedblright{}.) Say it once more.
\end{tcolorbox}
